\begin{frame}[label=q5d2]{5d.2 · Funciones: sum, lit y round}
  \small
  \pregunta{¿Qué participación tiene cada periodo sobre el total de personas reportadas?}
  \entrada{93\,831 filas. \texttt{periodo}, \texttt{cantidad}; denominador nacional 121\,427.}
  \flujo{Obtener total nacional\\\texttt{agg(sum)}}{\texttt{groupBy(periodo)}\\Suma por periodo}{\texttt{lit(totalNacional)}\\\texttt{round(100 * parte / total)}}{3 participaciones\\Suman 100\%}
  \begin{columns}[T,onlytextwidth]
    \column{0.55\textwidth}
    \textbf{Código Scala original (fragmento)}
    \codigo{sql/05d_2_participacion_periodo}{16}{19}
    \column{0.42\textwidth}
    \textbf{Salida}\\[5pt]
    {\footnotesize
\begin{tabular}{@{}lll@{}}
\toprule
Periodo & Personas & \% total\\
\midrule
2019 & 26\,331 & 21.68\\
2020--2021 & 38\,622 & 31.81\\
2022--2024 & 56\,474 & 46.51\\
\bottomrule
\end{tabular}
    }
  \end{columns}
  \respuesta{2022--2024 representa 46.51\% del acumulado. Cubre tres años, frente a dos y uno de los otros periodos; ese peso no demuestra mayor intensidad anual.}
  \origen{sql}{05d_2_participacion_periodo}
\end{frame}
